\section{BPMN-диаграмма процесса}
\headingtextsep

В качестве ключевого бизнес-процесса для BPMN-моделирования выбран процесс обработки задачи от постановки до завершения. Такой выбор обусловлен тем, что именно задача является центральной сущностью системы «Контур», а ее жизненный цикл объединяет действия координатора, исполнителя и самой системы.

\textheadingsep
\subsection{Участники процесса}
\headingtextsep

В BPMN-диаграмме используются три дорожки: «Координатор», «Исполнитель» и «Система «Контур»». Координатор отвечает за постановку и приемку результата. Исполнитель выполняет практическую работу по задаче и передает итог на проверку. Система обеспечивает валидацию данных, контроль зависимостей, отправку уведомлений, отслеживание сроков и фиксацию событий аудита. Сама диаграмма приведена в приложении Б на рисунке \ref{fig:appendix-bpmn-task-lifecycle}.

\textheadingsep
\subsection{Структура процесса}
\headingtextsep

Процесс начинается стартовым событием постановки новой задачи координатором. Далее выполняется последовательность действий и решений:

\begin{itemize}
\item координатор вводит атрибуты задачи и инициирует ее создание;
\item система выполняет валидацию данных и через эксклюзивный шлюз выбирает одну из ветвей: [данные корректны] или [обнаружены ошибки заполнения];
\item при корректных данных задача сохраняется, а система фиксирует событие создания в журнале аудита;
\item затем через второй эксклюзивный шлюз определяется, есть ли у задачи блокирующие зависимости: [блокировок нет] или [задача заблокирована];
\item при наличии блокировки запускается ожидание с промежуточным таймерным событием, после которого система повторно проверяет состояние зависимостей;
\item если блокировки отсутствуют, процесс проходит через параллельный шлюз: в одной ветви система отправляет уведомление исполнителю, а в другой начинает контроль срока выполнения и формирование напоминаний;
\item исполнитель принимает задачу в работу, выполняет необходимые действия и передает результат на проверку через промежуточное событие сообщения;
\item координатор рассматривает результат и через третий эксклюзивный шлюз принимает решение: [результат принят] или [требуется доработка];
\item при необходимости доработки задача возвращается исполнителю и цикл выполнения повторяется;
\item при успешной приемке система переводит задачу в завершенное состояние, записывает итоговое событие аудита и завершает процесс конечным событием.
\end{itemize}

В выбранном процессе отражены как штатные, так и исключительные ситуации. К исключениям относятся ошибки первичной валидации, ожидание снятия блокирующих зависимостей, просрочка срока выполнения с отправкой напоминания и возврат результата на доработку после проверки координатором. Такая структура позволяет на BPMN-диаграмме показать не только линейный маршрут выполнения задачи, но и реальные точки принятия решений, параллельную работу системы и возвраты по циклу управления.